\section*{Appendix A. Laporan Implementasi Python Prototype}

Lampiran ini merangkum implementasi Python untuk prototype AMETM T--S fuzzy secara ringkas: arsitektur proyek, fungsi inti, serta validasi berbasis unit test dan pemeriksaan manual.

\subsection*{A.1 Struktur Proyek Implementasi}

\begin{lstlisting}[
    basicstyle=\ttfamily\footnotesize,
    numbers=none,
    breaklines=true,
    frame=single,
    caption={Struktur direktori utama implementasi Python tugas 4}
]
folder_proyek/
|-- pyproject.toml
|-- src/
|   `-- ametm_ts/
|       |-- cli.py
|       |-- fuzzy.py
|       |-- ametm.py
|       |-- simulation.py
|       |-- metrics.py
|       `-- wind_turbine.py
|-- tests/
|   |-- test_fuzzy.py
|   |-- test_ametm.py
|   |-- test_data_and_simulation.py
|   `-- test_metrics.py
|-- data/
|   |-- synthetic_wind_turbine_case.csv
|   `-- simulation_results.csv
`-- plots/
    `-- simulation_overview.png
\end{lstlisting}

\subsection*{A.2 Alur Eksekusi Pipeline}

Pipeline pada \texttt{src/ametm\_ts/cli.py} berjalan dalam urutan:
\begin{enumerate}
    \item Membuat data sintetik jika file input belum tersedia.
    \item Menjalankan simulasi plant + kontrol event-triggered adaptif.
    \item Menyimpan hasil simulasi ke CSV.
    \item Membuat visualisasi ringkasan.
    \item Menghitung metrik evaluasi.
\end{enumerate}

\subsection*{A.3 Snippet Fungsi Utama}

\textbf{1) Bobot fuzzy dua aturan (\texttt{fuzzy.py})}

\begin{lstlisting}[
    language=Python,
    basicstyle=\ttfamily\small,
    keywordstyle=\color{blue}\bfseries,
    stringstyle=\color{orange},
    commentstyle=\color{gray}\itshape,
    breaklines=true,
    frame=single,
    caption={compute\_two\_rule\_weights(...)}
]
def compute_two_rule_weights(rotor_speed, nominal_speed=12.0, spread=4.0):
    if spread <= 0:
        raise ValueError("spread must be positive")

    left = nominal_speed - spread / 2.0
    right = nominal_speed + spread / 2.0

    if rotor_speed <= left:
        low = 1.0
    elif rotor_speed >= right:
        low = 0.0
    else:
        low = (right - rotor_speed) / (right - left)

    high = 1.0 - low
    weights = np.array([low, high], dtype=float)
    weights /= np.sum(weights)
    return weights
\end{lstlisting}

\textbf{2) Ambang adaptif AMETM (\texttt{ametm.py})}

\begin{lstlisting}[
    language=Python,
    basicstyle=\ttfamily\small,
    keywordstyle=\color{blue}\bfseries,
    stringstyle=\color{orange},
    commentstyle=\color{gray}\itshape,
    breaklines=true,
    frame=single,
    caption={adaptive\_threshold(...) dan should\_release(...)}
]
def adaptive_threshold(self, ack_attack_flag: int) -> float:
    history_mean = float(np.mean(self._history))
    attack_term = self.attack_scale if int(ack_attack_flag) == 1 else 0.0
    return self.base_threshold * (1.0 + attack_term + self.memory_weight * history_mean)

def should_release(self, state_error, ack_attack_flag: int) -> bool:
    threshold = self.adaptive_threshold(ack_attack_flag)
    error_energy = float(state_error.T @ state_error)
    return error_energy > threshold
\end{lstlisting}

\textbf{3) Logika event-triggered + ZOH saat DoS (\texttt{simulation.py})}

\begin{lstlisting}[
    language=Python,
    basicstyle=\ttfamily\small,
    keywordstyle=\color{blue}\bfseries,
    stringstyle=\color{orange},
    commentstyle=\color{gray}\itshape,
    breaklines=true,
    frame=single,
    caption={Ekstrak loop utama run\_simulation(...)}
]
release_attempt = trigger.should_release(
    state_error=(state_before - last_released_state),
    ack_attack_flag=ack_flag,
)

dos_active = dos_flag == 1
release_success = bool(release_attempt and not dos_active)

if release_success:
    held_control = control_candidate
    last_released_state = state_before.copy()

trigger.update_memory(release_success)
state_after = plant.step(float(held_control), disturbance, weights)
\end{lstlisting}

\textbf{4) Ringkasan metrik evaluasi (\texttt{metrics.py})}

\begin{lstlisting}[
    language=Python,
    basicstyle=\ttfamily\small,
    keywordstyle=\color{blue}\bfseries,
    stringstyle=\color{orange},
    commentstyle=\color{gray}\itshape,
    breaklines=true,
    frame=single,
    caption={summarize\_metrics(...)}
]
state_norm = np.sqrt(results["x1"] ** 2 + results["x2"] ** 2)
release_ratio = float(results["release_success"].mean())
bounded_ratio = float((state_norm <= state_bound).mean())

return {
    "release_ratio": release_ratio,
    "communication_saving": 1.0 - release_ratio,
    "bounded_ratio": bounded_ratio,
    "rms_state": float(np.sqrt(np.mean(state_norm**2))),
    "control_effort_l1": float(np.mean(np.abs(results["u_applied"]))),
}
\end{lstlisting}

\textbf{5) Orkestrasi CLI (\texttt{cli.py})}

\begin{lstlisting}[
    language=Python,
    basicstyle=\ttfamily\small,
    keywordstyle=\color{blue}\bfseries,
    stringstyle=\color{orange},
    commentstyle=\color{gray}\itshape,
    breaklines=true,
    frame=single,
    caption={Ekstrak run\_pipeline(...)}
]
if not data_path.exists():
    synthetic = generate_synthetic_case(num_steps=num_steps, dt=dt, seed=seed)
    save_case_csv(synthetic, data_path)

case_data = load_case_csv(data_path)
results = run_simulation(case_data, plant=plant, trigger=trigger, gains=gains)
save_case_csv(results, result_path)
plot_simulation_results(results, plot_path)

return summarize_metrics(results)
\end{lstlisting}

\subsection*{A.4 Validasi Berbasis Code (Unit Test)}

Pengujian otomatis dilakukan dengan \texttt{unittest} pada empat berkas:
\begin{itemize}[leftmargin=*]
    \item \texttt{tests/test\_fuzzy.py}: validasi normalisasi bobot fuzzy dan batas ekstrem.
    \item \texttt{tests/test\_ametm.py}: validasi kenaikan threshold pada status serangan dan keputusan release.
    \item \texttt{tests/test\_data\_and\_simulation.py}: validasi skema data input dan konsistensi output simulasi.
    \item \texttt{tests/test\_metrics.py}: validasi rentang metrik dan ketersediaan key output.
\end{itemize}

Perintah uji:
\begin{lstlisting}[
    basicstyle=\ttfamily\small,
    numbers=none,
    breaklines=true,
    frame=single,
    caption={Eksekusi unit test}
]
python -m unittest discover -s tests -t . -p "test_*.py"
\end{lstlisting}

Ringkasan hasil eksekusi:
\begin{lstlisting}[
    basicstyle=\ttfamily\small,
    numbers=none,
    breaklines=true,
    frame=single,
    caption={Output unit test}
]
.......
----------------------------------------------------------------------
Ran 7 tests in 0.006s

OK
\end{lstlisting}

\subsection*{A.5 Validasi Manual Hasil Simulasi}

Selain unit test, dilakukan verifikasi manual terhadap metrik dari file
\texttt{data/simulation\_results.csv}. Konsistensi utama dicek dengan relasi:
\[
\texttt{communication\_saving} = 1 - \texttt{release\_ratio}
\]

Hasil verifikasi manual sesuai output pipeline:

\begin{center}
\begin{tabular}{l r}
\toprule
\textbf{Metrik} & \textbf{Nilai} \\
\midrule
release\_ratio & 0.010000 \\
communication\_saving & 0.990000 \\
bounded\_ratio & 1.000000 \\
rms\_state & 0.046057 \\
control\_effort\_l1 & 0.025128 \\
\bottomrule
\end{tabular}
\end{center}

Pemeriksaan manual tambahan pada membership dua aturan menunjukkan normalisasi konsisten ($w_{low}+w_{high}=1$):

\begin{center}
\begin{tabular}{c c c c}
\toprule
\textbf{Kecepatan Rotor} & $w_{low}$ & $w_{high}$ & Jumlah \\
\midrule
8.0  & 1.0000 & 0.0000 & 1.0000 \\
12.0 & 0.5000 & 0.5000 & 1.0000 \\
16.0 & 0.0000 & 1.0000 & 1.0000 \\
\bottomrule
\end{tabular}
\end{center}

\subsection*{A.6 Ringkasan Lampiran}

Implementasi Python telah dipisahkan per modul fungsional (fuzzy, AMETM,
simulasi, metrik, dan CLI), dilengkapi validasi unit test, serta diverifikasi
manual pada metrik kunci. Struktur ini membuat prototype mudah diuji ulang,
dikembangkan, dan dipetakan kembali ke konsep teoretis yang dibahas pada
review utama.
